% !PRINTED 347
\weberpage{347}
%<<CONT
so $P_0$ is no longer transitive, that is, after adjunction of $x_0$
the equation for the remaining eight roots splits into several factors. If the
group $P$ is transitive at all, then $x_0$ can here be any one of the roots;
for if $\pi$ is a permutation of $P$ by which $x_0$ passes into $x_1$,
then by $P_1 = \pi^{-1} P_0 \pi$ the root $x_1$ remains unchanged, and if $P_0$
is intransitive for eight roots, then $P_1$ must also be intransitive for eight roots,
since the permutations of $P_1$ differ from those of $P_0$ only in the designation of the
roots (Bd. I, §.\,154, 6.).

The simply transitive ones among the linear groups have a more special character
than the group of the triple equations.

In §.\,78 we became acquainted with a decomposition of the general linear congruence
group for the modulus 3, from which we can survey these relations.

The \gesperrt{Abelian} group $Q$, which consists of the substitutions

\weq{9}{\xi' \equiv \xi + \alpha, \quad \eta' \equiv \eta + \beta \quad (\mathrm{mod}\ 3)}

is certainly only simply transitive; for if a root is to remain unchanged
by (9), then we must have $\alpha \equiv 0$, $\beta \equiv 0$, and all roots remain
unchanged. Every root is rationally expressible in terms of every other.

If we now seek out in the group $P$ the group $P_0$ of the substitutions by which
$(0,0)$ remains unchanged, then $P_0$ contains only homogeneous substitutions and is
therefore the greatest common divisor of $P$ and the homogeneous congruence group $S$,
and we have $P = P_0 Q$.

If we form for $S$ the composition series (§.\,78):

\[
S,\ S_1,\ S_2,\ S_3,\ S_4,
\]

then $S_4 = \begin{pmatrix} 1, & 0 \\ 0, & 1 \end{pmatrix}, \begin{pmatrix} -1, & 0 \\ 0, & -1 \end{pmatrix}$,
and by $S_4$ the roots $(1,1)$, $(2,2)$ merely pass into one another. By $S_4 Q$
the triple $(0,0)$ $(1,1)$ $(2,2)$ can indeed undergo every arrangement of its roots.
But it can pass over only into the two other triples $(1,0)$ $(2,1)$ $(0,2)$;
$(0,1)$ $(1,2)$ $(2,0)$. $S_4 Q$ is therefore only \gesperrt{simply transitive} and
is moreover \gesperrt{imprimitive}.

By the group $S_3$, which arises from $S_4$ by the addition of the
substitution

\[
\begin{pmatrix} -1, & 1 \\ 1, & 1 \end{pmatrix}
\]

%CONT>>
